% !PRINTED 15
\weberpage{15}

Die Construction eines Schnittes in der Reihe der Zahlen liefert
stets den Beweis für die Existenz einer Zahl, die bestimmten
Anforderungen genügt. So erhält man einen Schnitt $(A,B)$, wenn man
alle und nur die Zahlen, deren Quadrat kleiner als eine bestimmte Zahl
$\alpha$ ist, in $A$ aufnimmt; diesem Schnitt entspricht eine bestimmte
Zahl, deren Quadrat gleich $\alpha$ ist und die mit $\sqrt\alpha$
bezeichnet wird, und dadurch wird die Existenz der Quadratwurzeln
nachgewiesen.

Auf die Schnitte lassen sich auch die von G.\ \gesperrt{Cantor} zur
Definition der Irrationalzahlen eingeführten Zahlenreihen zurückführen.\setcounter{footnote}{0}\footnote{
\gesperrt{Cantor}, Ueber die Ausdehnung eines Satzes aus der Theorie der
trigonometrischen Reihen. Mathematische Annalen, Bd.~5 (1872); vergl.
auch \gesperrt{Heine}, Elemente der Functionenlehre. Journal für
Mathematik, Bd.~74 (1872).}

Nach \gesperrt{Cantor} ist unter einer Zahlenreihe irgend ein
unbegrenztes, in bestimmter Weise geordnetes System von Zahlen zu
verstehen:
\[
 S=x_1,x_2,x_3,x_4,\ldots
\]
von der Beschaffenheit, dass es eine bestimmte Zahl $g$ giebt, unter
die keine der Zahlen $S$ heruntersinkt, und dass, wenn $\delta$ eine
beliebig gewählte Zahl ist, und $x_n,x_m$ verschiedene Elemente aus
$S$ sind, $x_n-x_m$ oder $x_m-x_n$ immer kleiner bleibt als $\delta$,
wenn $m$ und $n$ eine hinlänglich grosse Zahl überschritten haben.

Es giebt immer Zahlen, die von den Zahlen einer solchen Reihe nicht
überschritten werden; denn hat man $\delta$ gewählt und $n$ passend
bestimmt, so überschreitet $x_m$, welchen Werth auch $m$ haben mag,
niemals die grösste der Zahlen $x_1,x_2,\ldots,x_n,x_n+\delta$. Man
erhält nun einen Schnitt $(A,B)$, wenn man die Zahlen nach $B$ wirft,
die, wenn $n$ einen hinlänglich hohen Werth hat, von keinem $x_n$ mehr
überschritten werden, und alle anderen Zahlen (die also von unendlich
vielen $x_n$ überschritten werden) nach $A$. Wird dieser Schnitt durch
die Zahl $\alpha$ erzeugt, so giebt es, wie klein auch $\varepsilon$
sei, immer unendlich viele Zahlen $x_n$ zwischen
$\alpha-\varepsilon$ und $\alpha$, und man kann sagen, dass diese durch
$S$ vollkommen bestimmte Zahl $\alpha$ durch die Zahlenreihe $S$
erzeugt wird. Nach \gesperrt{Cantor} ist die Zahlenreihe $S$ geradezu
die Definition der Zahl $\alpha$. Selbstverständlich kann eine und dieselbe
%CONT>>
